\chapter{2004 Nobel Prize in Economics}
\textbf{Laureates: }Finn Kydland (Norway), Edward Prescott (USA)

\section*{Reason for Award}
For their contributions to dynamic macroeconomics: the time consistency of
economic policy and the driving forces behind business cycles

\section{What They Did}
Kydland and Prescott made two groundbreaking contributions that reshaped
macroeconomic theory, policy analysis, and the methodology of quantitative economics.
Kydland, born in Norway in 1943, and Prescott, born in Maine in 1940, developed
theoretical frameworks that have become foundational to modern macroeconomics.

First, they formalized the time inconsistency problem in economic policy,
published in their seminal 1977 paper, "Rules Rather Than Discretion: The
Inconsistency of Optimal Plans." They showed that policymakers facing
discretionary choice have an incentive to deviate from previously announced
optimal plans once economic agents have formed expectations based on those
plans. This creates a fundamental tension between what is optimal ex ante
and what is optimal ex post.

In the context of monetary policy, Kydland and Prescott demonstrated the
inflation bias that arises under discretion. A central bank may announce a
low inflation target to anchor expectations. Once workers and firms set wages
and prices based on that target, the central bank has an incentive to create
surprise inflation to boost output above its natural rate. Anticipating this
incentive, rational agents build higher inflation into their wage and price
settings from the start, resulting in higher inflation without any gain in
output. This insight showed that discretion can lead to worse outcomes than
rules.

Second, Prescott developed the Real Business Cycle (RBC) theory, which
identified technology shocks as the primary driver of business fluctuations.
Building on earlier work by Robert Lucas and others, Prescott showed that
fluctuations in output, employment, and investment could arise as optimal
responses to real shocks--primarily changes in total factor productivity--within
a general equilibrium framework. This challenged the Keynesian view that
business cycles are primarily driven by demand shocks.

Prescott developed quantitative business cycle theory, calibrating model
parameters to match observed moments of the data and simulating the model's
implications. He showed that RBC models could reproduce the observed
autocorrelations, volatilities, and cross-correlations of macroeconomic
variables with remarkable accuracy. This quantitative approach became the
standard for macroeconomic model building.

Their work also addressed the optimal taxation problem, showing how tax
policies affect incentives for work, saving, and investment. The Ramsey
taxation problem they analyzed has important implications for the design
of optimal tax systems.

\section{Impact on Economic Thinking}
Kydland and Prescott's time inconsistency insight became central to monetary
policy design. Central bank independence, adopted by numerous countries
including Canada, New Zealand, the United Kingdom, and eventually the European
Central Bank, directly responds to their analysis. Inflation targeting regimes,
which emerged in the 1990s, incorporate commitment mechanisms to anchor
expectations and enhance credibility.

Their work showed that credible commitment is essential for effective monetary
policy. The inflation bias under discretion provides a theoretical justification
for central bank independence and rules-based policy frameworks. This insight
has influenced institutional design of monetary policy across the world.

Prescott's RBC framework dominated academic macroeconomics through the 1980s
and 1990s. It introduced rigorous quantitative methods, emphasizing calibration,
simulation, and the comparison of model-implied moments with data. The New
Keynesian synthesis incorporated RBC foundations while adding nominal
rigidities--sticky prices and wages--and demand shocks, creating the dominant
Dynamic Stochastic General Equilibrium (DSGE) framework used by central banks today.

The quantitative macroeconomics approach that Prescott championed transformed
the profession. Macroeconomists now routinely calibrate models, simulate
counterfactual policies, and compare model implications with empirical moments.
This methodological approach has raised the standards of empirical work in
macroeconomics.

The time inconsistency problem they identified has applications beyond monetary
policy, including fiscal policy, regulatory policy, and international
cooperation. Their insights about commitment devices have informed debates
about constitutional rules, fiscal rules, and institutional design.

\section{Modern Perspectives Today}
DSGE models incorporating RBC foundations remain the standard analytical
framework at central banks worldwide. The Federal Reserve, the European
Central Bank, and other major monetary institutions rely on DSGE models for
policy analysis, forecasting, and the evaluation of alternative policy rules.
The COVID-19 pandemic renewed interest in DSGE models for policy analysis.

Fiscal policy debates reference optimal tax smoothing and the constraints on
discretionary policy. Discussions about fiscal rules, balanced budget
amendments, and independent fiscal councils reflect Kydland and Prescott's
insights about the value of commitment devices. Their work has influenced
debates about fiscal sustainability and the design of fiscal rules.

Real business cycle theory informs supply-side policy evaluations,
infrastructure analysis, and productivity research. The emphasis on
microfoundations, general equilibrium, and intertemporal optimization continues
to shape macroeconomic modeling. Quantitative methods, with their emphasis on
calibration and simulation, remain central to the discipline.

Their work continues to influence debates about the appropriate role of
government in the economy and the design of institutional arrangements for
economic management. The insights about time inconsistency and commitment
remain relevant for understanding policy challenges in an uncertain world.
